\documentclass[11pt]{article}
\usepackage[margin=1in]{geometry}
\usepackage[T1]{fontenc}
\usepackage{lmodern}
\usepackage{booktabs}
\usepackage{multirow}
\usepackage{xurl}
\usepackage[hidelinks]{hyperref}
\usepackage{amsmath}

\title{Verifying Attested Packages Without the Network:\\
A Paired Rebuild Study of npm and PyPI}
\author{[Author name]\\ \small [Affiliation] \quad \texttt{[email]}}
\date{November 2026}

\begin{document}
\maketitle

\begin{abstract}
Build provenance on npm and PyPI says which source commit and workflow
produced a package, and a verifier can check that claim by rebuilding the
package and comparing the result with what the registry serves. A
rebuild that downloads dependencies from the network, though, can be
steered by whoever controls those downloads. We measure how many attested
releases still verify when the rebuild runs with the network off and every
dependency must come from a hash-pinned lock in the attested commit. For
600 npm and 600 PyPI releases drawn at random from the deps.dev snapshot of
9 November 2026, each release was rebuilt three ways with the same
recipe: network off, network on, and network off with a lock inferred at
publish time. On npm, 22.3\% of releases verify with the network off and
30.2\% with it, a paired difference of 7.8 percentage points (95\% CI 5.7
to 10.0; exact McNemar $p = 1.4\times10^{-14}$). No release verified
offline that failed online. On PyPI, no release verifies offline because
none commits a build lock, while 71.2\% verify online and 69.5\% verify
offline from an inferred lock. Almost all of the gap is a missing or
incomplete lock, not a build that needs the network.
\end{abstract}

\section{Introduction}

Attacks on open-source supply chains often go through the path from source
to published artifact: a compromised build machine, a stolen publishing
token, or a release step that adds code not in the repository
\cite{ohm2020,ladisa2023}. Build provenance answers part of this. npm and
PyPI now accept Sigstore-signed statements \cite{newman2022} in the SLSA
format \cite{slsa,npmprov,pep740} that name the repository, commit and
workflow behind a release. A signature shows who built the package. It
does not show that the package matches the commit. For that, a verifier
has to rebuild from the commit and compare
\cite{lamb2022,ossrebuild}.

Two recent studies measure how often such rebuilds succeed across
ecosystems \cite{solarin2026,dietrich2026}. Neither reports whether its
rebuilds had network access. The question matters for security. A rebuild
that resolves dependencies from the registry trusts the registry and the
network at verification time, and can be made to agree with a tampered
artifact. A rebuild that runs with the network off, installing only
dependencies whose hashes are fixed by a lock in the attested commit,
trusts only what the provenance names.

We ask how much verifiability that stricter setting costs.

\begin{itemize}
\item \textbf{Q1.} For attested npm and PyPI releases, how many verify
with the network off, how many with it on, and how large is the paired
difference?
\item \textbf{Q2.} Where do releases fail without the network: missing
lock, lock without hashes, a lock entry that needs the network, the build,
or the comparison?
\end{itemize}

The study plan, including the sample, arms, outcomes and test, was fixed
and committed before any rebuild ran; deviations are logged with dates in
the repository. A third question, on what the comparison rules used by
rebuild verifiers let through, is under coordinated disclosure with the
maintainers concerned and will be added in a later version of this paper.

\section{Method}

\subsection{Sample}

The frame is every npm and PyPI package in the deps.dev public dataset
\cite{depsdev} (snapshot 9 November 2026) with a verified attestation that
names a commit. For each package we take its most recent such release
published on or before 31 October 2026. Packages are ordered by
\texttt{FARM\_FINGERPRINT} of a fixed seed string and the package name,
so the draw can be repeated exactly. We take the first 600 npm packages.
deps.dev records about one attestation per PyPI version, usually a
wheel's, so PyPI candidates were screened in draw order through PyPI's
Integrity API, keeping releases with exactly one sdist that has its own
provenance, until 600 were kept. Of 629 screened, 27 had no sdist and 2
were missing from PyPI; none had an sdist without provenance. At 600
releases per ecosystem, a 95\% Wilson interval for any single rate is at
most $\pm4.0$ percentage points wide.

\subsection{Rebuilds}

Each release was rebuilt in three arms by ClearHash, an open-source
verifier, with the same tool version, container image and recipe. Only
the source of dependencies differed.

\begin{description}
\item[Arm H (network off).] The shipped verifier. Dependencies come only
from a lock committed at the attested commit (npm, pnpm or Yarn locks with
SHA-256 or stronger hashes; for PyPI, a \texttt{pylock.toml}
\cite{pylock} covering the build backend). The host downloads and checks
each archive; the build container has no network.
\item[Arm N (network on).] The container can reach the registry. npm runs
\texttt{npm ci} with a lock, or \texttt{npm install --before} the publish
time without one; pip resolves with \texttt{PIP\_UPLOADED\_PRIOR\_TO} set
to the publish time.
\item[Arm I (inferred lock, network off).] When the commit has no lock arm
H accepts, a networked container writes a hash-pinned lock as of the
publish time, and the build then runs as in arm H. Arm I separates
``the project committed no usable lock'' from ``the build needs the
network.''
\end{description}

Sources are fetched at depth one, as \texttt{actions/checkout} does by
default, and the checked-out commit must equal the attested one. Each rebuild
has a 30-minute limit. The build host is Linux on ARM64.

\subsection{Outcomes and analysis}

A release \emph{verifies} in an arm when ClearHash's normalized file-tree
comparison (T1) finds the rebuilt package equal to the registry's. T1
compares the regular files in both packages, their contents and their
owner-execute bit, not archive timestamps or order. At the frozen
version it also drops registry-added \texttt{package.json} fields and
compares that file with its keys sorted, ignores npm's root
\texttt{.npmignore} and an sdist's \texttt{PKG-INFO},
\texttt{*.egg-info} and \texttt{*.dist-info} files, and skips symbolic
links. We also report T0, a strict comparison without any of those
rules, in which a symbolic link counts with its target. Every failed rebuild is assigned the first
stage it failed at, and a reason within the stage from its error text,
using a fixed mapping committed with the analysis code.

The confirmatory test is an exact McNemar test \cite{mcnemar1947} of arm H
against arm N per ecosystem, Holm-corrected \cite{holm1979} across
ecosystems at $\alpha = 0.05$. The plan made an ecosystem with fewer
than 10 arm-H matches descriptive only. Single rates carry Wilson
intervals \cite{wilson1927} and paired differences Newcombe's
interval for paired proportions \cite{newcombe1998}. Because some
repositories publish several sampled packages (one npm monorepo
contributes 15), we repeat the primary estimate with a repository-level
cluster bootstrap and with one release per repository.

\subsection{Run and reruns}

The main run produced all 3,600 records with the tool frozen at a tagged
commit. Reruns followed rules written before any match rate was examined:
two records that failed on Docker errors were rerun; 73 releases that
stopped at the source stage were retried once and none recovered; and a
second tagged version, fixing a hash check for pinned Yarn releases and
isolating Git from the host's configuration, reran the 11 releases those
two bugs affected. No release is missing from the analysis.

\section{Results}

\subsection{Q1: verification with and without the network}

\begin{table}[t]
\centering
\caption{Releases that verify (T1), by arm. Brackets: 95\% Wilson
interval.}
\label{tab:rates}
\begin{tabular}{lrrr}
\toprule
& Arm H (off) & Arm N (on) & Arm I (inferred) \\
\midrule
npm ($n=600$) & 134 (22.3\%) & 181 (30.2\%) & 162 (27.0\%) \\
& [19.2, 25.8] & [26.6, 34.0] & [23.6, 30.7] \\
PyPI ($n=600$) & 0 (0.0\%) & 427 (71.2\%) & 417 (69.5\%) \\
& [0.0, 0.6] & [67.4, 74.6] & [65.7, 73.1] \\
\midrule
npm, strict (T0) & 128 & 176 & 156 \\
PyPI, strict (T0) & 0 & 419 & 409 \\
\bottomrule
\end{tabular}
\end{table}

Table~\ref{tab:rates} gives the rates. On npm, 134 releases verify in
both arm H and arm N, 47 only in arm N, none only in arm H, and 419 in
neither. The paired difference is 7.8 percentage points (95\% CI 5.7 to
10.0), and the exact McNemar test gives $p = 1.4\times10^{-14}$; with npm
the only ecosystem in the test family, the Holm-adjusted $p$ is the same.
Clustering does not change this: the repository bootstrap gives 5.8 to
10.1 over 569 repositories, and counting one release per repository gives
8.3 points (6.0 to 10.6).

PyPI met the degenerate-case rule. No release verified in arm H, because
none of the 600 commits a \texttt{pylock.toml}. With the network, 427
verify, and with an inferred lock and the network off, 417 do. PyPI builds
therefore work offline almost as well as online once the build backend is
pinned; the obstacle is that projects do not commit that pin.

The secondary contrasts show the same thing on npm. An inferred lock
recovers 28 of the 47 releases that verify only with the network (arm I
minus arm H: 4.7 points, 3.0 to 6.4), and arm N still verifies 19 more
than arm I (3.2 points, 1.8 to 4.6). On PyPI, arm N exceeds arm I by 10
releases (1.7 points, 0.6 to 2.7). Where arms N and I both verify a
release, they produce the same tree.

Strict comparison changes little. T0 is within 6 releases of T1 in every
npm arm and within 8 in every PyPI arm.

\subsection{Q2: where offline rebuilds fail}

\begin{table}[t]
\centering
\caption{First failing stage and reason, by arm. Reasons come from a
fixed mapping of error text; ``other'' is anything it does not match.}
\label{tab:stages}
\small
\begin{tabular}{llrrrrrr}
\toprule
& & \multicolumn{3}{c}{npm} & \multicolumn{3}{c}{PyPI} \\
\cmidrule(lr){3-5}\cmidrule(lr){6-8}
Stage & Reason & H & N & I & H & N & I \\
\midrule
Attestation & workflow in another repository & 32 & 32 & 32 & 0 & 0 & 0 \\
Fetch & artifact missing from registry & 1 & 1 & 1 & 0 & 0 & 0 \\
Source & attested commit unavailable & 22 & 22 & 22 & 42 & 42 & 42 \\
\midrule
\multirow{5}{*}{Lock}
& no committed lock & 130 & -- & -- & 558 & -- & -- \\
& lock without usable hashes & 26 & -- & 28 & 0 & -- & 0 \\
& lock entry needs the network & 13 & -- & 13 & 0 & -- & 0 \\
& unsupported lock version & 6 & -- & 6 & 0 & -- & 0 \\
& lock inference failed / other & 2 & -- & 31 & 0 & -- & 24 \\
\midrule
Dependencies & fetch or integrity & 2 & 1 & 2 & 0 & 0 & 0 \\
\midrule
\multirow{6}{*}{Build}
& version not in source tree & 89 & 156 & 129 & 0 & 0 & 0 \\
& TypeScript compile failed & 38 & 50 & 40 & 0 & 0 & 0 \\
& build tried to use the network & 10 & 0 & 11 & 0 & 0 & 0 \\
& engine version mismatch & 6 & 9 & 7 & 0 & 0 & 0 \\
& build backend failed & 0 & 0 & 0 & 0 & 54 & 43 \\
& timeout or other & 31 & 76 & 55 & 0 & 0 & 0 \\
\midrule
Artifact & unexpected archive layout & 0 & 0 & 0 & 0 & 38 & 37 \\
\midrule
\multirow{2}{*}{Compared}
& verifies & 134 & 181 & 162 & 0 & 427 & 417 \\
& differs & 58 & 72 & 61 & 0 & 39 & 37 \\
\bottomrule
\end{tabular}
\end{table}

Table~\ref{tab:stages} gives every release's first failing stage. Four
results stand out.

\emph{The lock is the main barrier.} Of the 47 npm releases that verify
only with the network, 45 stop at the lock stage in arm H: 36 have no
committed lock, 4 lock dependencies without usable hashes, 4 lock an entry
the host cannot fetch from the registry (a Git or arbitrary-URL source),
and 1 uses an unsupported lock version. One more fails because its build
script downloads at build time, and one builds offline but produces a
different tree. Across all npm releases, 130 commit no lock at all.

\emph{Few builds need the network beyond their dependencies.} Only 10
npm rebuilds failed in arm H because the build itself reached for the
network. With the network, only one of those 10 verifies; four build but
differ from the registry, and five fail to build.

\emph{Many failures are independent of the network.} The largest npm
build failure, in every arm, is that the package's published version does
not appear in any \texttt{package.json} at the attested commit, because
the publish workflow sets the version. A rebuild from the commit cannot
reproduce such a package without replaying that step. TypeScript and
Node-engine failures are similar: the recipe is not the publisher's
workflow.

\emph{Provenance is not always checkable later.} For 64 releases (22
npm, 42 PyPI) the attested commit can no longer be fetched, mostly because
the repository has been deleted or made private. For 32 npm releases the
signing workflow lives in a different repository from the attested
source (a shared release workflow), which ClearHash rejects because it
cannot tie the signer to the source.

\section{Limitations}

Absolute rates are lower bounds. The recipe is a generic one, not each
publisher's workflow, and the build host is Linux on ARM64, so packages
that build only elsewhere fail in every arm. The paired difference is not
affected by either, because all three arms share them. The sample covers
releases with verified provenance only, from one snapshot, and the frame
reflects what deps.dev recorded. Sources are fetched at depth one, so
builds that read Git history see a single commit in every arm. The
retry of source failures recovered none, so they are not transient, but
we did not establish the reason for every one. Finally, T1 depends on
ClearHash's comparison rules at the frozen version; T0 shows the strict
counts.

\section{Related work}

Solarin et al.\ \cite{solarin2026} measure how many artifacts in four
ecosystems can be verified from registry metadata at tiered equivalence
levels, and Dietrich et al.\ \cite{dietrich2026} compare alternatively
built PyPI wheels. OSS Rebuild \cite{ossrebuild} rebuilds packages from
several ecosystems and publishes attestations of the result. These works
report how often rebuilds succeed; this study isolates what network
access at verification time contributes, with the same releases rebuilt
both ways.

\section{Conclusion}

Verifying attested npm and PyPI releases with the network off costs
little in what the builds need and a lot in what projects commit. On npm,
offline verification covers about three quarters of what networked
verification does, and the difference is almost entirely missing or
incomplete locks. On PyPI, inferred locks show that builds would verify
offline at nearly the networked rate, if projects committed a lock for
their build backend. Registries that accept provenance could ask for that
lock at publish time.

\section*{Availability}

The study plan, sample, analysis code, and the per-release ledger are in
the ClearHash repository (\url{https://github.com/Builder106/clear-hash}),
with the tool versions used tagged \texttt{hv-freeze} and
\texttt{hv-freeze-2}.

\bibliographystyle{plain}
\bibliography{refs}
\end{document}
